\Chapter{4}

\Verse{1}
{
	\hJ{וְהָ}
	\hJ{\hlB{אָדָם}}
	\hJ{\hlC{יָדַע}}
	\hJ{אֶת חַוָּה אִשְׁתּוֹ וַתַּהַר וַתֵּלֶד אֶת קַיִן וַתֹּאמֶר קָנִיתִי אִישׁ אֶת יְהֹוָה}
}
{
	\eJ{
		And the \hlB{human} had \hlC{known} Eve, his woman,
		and she became pregnant and gave birth to Cain,
		and said, “I’ve created a man with YHWH.”
	}
}
{
	Then the woman has the first baby ever.

	And presumably she was surprised, because no one ever made a person before.

	And so naturally she goes hey ``I've created a man like YHWH.''

	And she names the kid Cain, which means something like ``maker'' or ``creator'' or ``smith.''

	And she says basically ``I've cain'ed a man like YHWH,'' because see line above.

	Also notice up in the first line how the noun \paleo{אדמ} (adm) comes before the verb \paleo{יָדַע} (yd`).

	And whenever that happens in biblical hebrew, it's the past tense plus the perfect aspect.

	And as you know, we didn't know how to read the past perfect in biblical hebrew until 1998.

	Anyways there's a point to all this grammar:

	They already had sex back in the garden.

	Which is what the author is telling us by putting the noun before the verb.

	That's why the English says ``had known'' and not ``knew.''

	Just make a note of that and let's keep moving.

}

\Verse{2}
{
	\hJ{וַתֹּסֶף לָלֶדֶת אֶת אָחִיו אֶת הָבֶל וַיְהִי הֶבֶל רֹעֵה צֹאן וְקַיִן הָיָה עֹבֵד אֲדָמָה}
}
{
	\eJ{
		And she went on to give birth to his brother, Abel.
		And Abel was a shepherd of flocks,
		and Cain was a worker of ground.
	}
}
{
	Apparently she got pregnant again while we were talking about grammar.%

	This one's name is HBL or HVL or Hevel but we write it Abel, because reasons.

	And it's the same word as

	\Def{HBL}[\emph{Hevel}, \heb{הבל}]{

		\aC{
			Traditionally rendered ``Abel'' in English, the term literally means
			\emph{breath}, \emph{vapor}, \emph{mist}, or something \emph{fleeting} and
			\emph{insubstantial}. It is the same word made famous by the book of Ecclesiastes.
			\heb{\hlC{הבל}ים} \heb{\hlC{הבל}}'' --- ``vanity of vanities,'' or more literally,
			``vapor of vapors.'' The traditional Ecclesiastical translation of ``vanity''
			means not the modern sense of narcissism, but rather the sense of the term
			meaning ``in vain.''
		}
		
		\aC{
			The name suggests something beautiful but short-lived: a puff of breath
			on a cold morning, a wisp of smoke, or a mist that disappears almost as soon
			as it appears. The irony is unmistakable. Hevel's life is as brief and
			fleeting as his name.
		}
	}

	\hphantom{And it's the same word as} in Qohelet%
	\footnote{
		A depressed drunktext traditionally attributed to king Solomon.
		A thousand wives will do that to you, according to some scholars.}
	where he keeps saying everything is this.

	The second baby's name means something like ``He's just temporary.''

}

\Verse{3}
{
	\hJ{וַיְהִי מִקֵּץ יָמִים וַיָּבֵא קַיִן מִפְּרִי הָאֲדָמָה מִנְחָה לַיהֹוָה}
}
{
	\eJ{
		And it was at the end of some days
		and Cain brought an offering to YHWH from
		the fruit of the ground.
	}
}
{
	So the first kid brings some vegetables to YHWH.
}

\Verse{4}
{
	\hJ{וְהֶבֶל הֵבִיא גַם הוּא מִבְּכֹרת צֹאנוֹ וּמֵחֶלְבֵהֶן וַיִּשַׁע יְהֹוָה אֶל הֶבֶל וְאֶל מִנְחָתוֹ}
}
{
	\eJ{
		And Abel brought, as well, from the firstborn of his flock and their fat.
		And YHWH paid attention to Abel and his offering.
	}
}
{
	And the other one brings some meat.
	
	And YHWH notices that one.
}

\Verse{5}
{
	\hJ{וְאֶל קַיִן וְאֶל מִנְחָתוֹ לֹא שָׁעָה וַיִּחַר לְקַיִן מְאֹד וַיִּפְּלוּ פָּנָיו}
}
{
	\eJ{
		And to Cain and his offering he did not pay attention.
		And Cain was very upset, and his face was fallen.
	}
}
{
	He didn't notice the vegetables.
}

\Verse{6}
{
	\hJ{וַיֹּאמֶר יְהֹוָה אֶל קָיִן לָמָּה חָרָה לָךְ וְלָמָּה נָפְלוּ פָנֶיךָ}
}
{
	\eJ{
		And YHWH said to Cain, “Why are you upset,\nl
		and why has your face fallen?”
	}
}
{
	But he notices Maker is sad.
}

\Verse{7}
{
    \hJ{הֲלוֹא אִם תֵּיטִיב שְׂאֵת וְאִם לֹא תֵיטִיב לַפֶּתַח חַטָּאת רֹבֵץ וְאֵלֶיךָ תְּשׁוּקָתוֹ וְאַתָּה תִּמְשׁל בּוֹ}
}
{
    \eJ{
        Is it not that if you do well you'll be raised,
        and if you don't do well then sin crouches at the threshold?
        And its desire will be for you.
        And you'll dominate it.
    }
}
{

    So YHWH says, well, whatever this is.

    \Def{Genesis 4:7}[]{

        \aC{
            One of the greatest misconceptions about Biblical Hebrew is the
            assumption that we understand it as well as Classical Greek or
            Latin. We do not. Hundreds of Hebrew words---and countless idioms,
            expressions, and passages---remain uncertain even to the world's
            leading scholars. Every translation must therefore make choices on
            the reader's behalf, sometimes explaining the uncertainty in a
            footnote, but more often leaving only the familiar refrain,
            ``Meaning of Heb. uncertain.'' Even when the meanings of Hebrew
            terms are well-known, it is often far from clear what an author
            was up to.
        }

        \begin{center}
        \begin{tabular}{p{0.22\linewidth}p{0.68\linewidth}}
        \hline
        \textbf{\aC{Hebrew}} & \textbf{\aC{Meaning}} \\
        \hline
        \hJ{הלא}      & \aC{surely? is it not?} \\
        \hJ{אם}       & \aC{if} \\
        \hJ{תיטיב}    & \aC{do well} \\
        \hJ{שאת}      & \aC{lifting; acceptance; forgiveness} \\
        \hJ{ואם}      & \aC{but if} \\
        \hJ{לא}       & \aC{not} \\
        \hJ{תיטיב}    & \aC{do well} \\
        \hJ{לפתח}     & \aC{at the doorway} \\
        \hJ{חטאת}     & \aC{sin; sin offering} \\
        \hJ{רבץ}      & \aC{crouching; lying down} \\
        \hJ{ואליך}    & \aC{toward you} \\
        \hJ{תשוקתו}   & \aC{desire; longing (\emph{3× in the Bible})} \\
        \hJ{ואתה}     & \aC{but you} \\
        \hJ{תמשל}     & \aC{rule; master} \\
        \hJ{בו}       & \aC{over it} \\
        \hline
        \end{tabular}
        \end{center}

    }

}

\Verse{8}
{
	\hJ{וַיֹּאמֶר קַיִן אֶל הֶבֶל אָחִיו וַיְהִי בִּהְיוֹתָם בַּשָּׂדֶה וַיָּקם קַיִן אֶל הֶבֶל אָחִיו וַיַּהַרְגֵהוּ}
}
{
	\eJ{
		And Cain said to his brother Abel.
		And it was while they were in the field,
		and Cain rose against Abel his brother
		and killed him.
	}
}
{
	And so Cain says.\footnote{Something happened here.}

	And they were in the field and he killed him.
}

\Verse{9}
{
	\hJ{וַיֹּאמֶר יְהֹוָה אֶל קַיִן אֵי הֶבֶל אָחִיךָ וַיֹּאמֶר לֹא יָדַעְתִּי הֲשֹׁמֵר אָחִי אָנֹכִי}
}
{
	\eJ{
		And YHWH said to Cain,
		“Where is Abel your brother?”
		And he said, “I don’t know.
		Am I my brother’s watchman?”
	}
}
{
	Gen 3:9, YHWH loses track of Adam.

	Gen 4:9, YHWH loses track of Abel.
}

\Verse{10}
{
	\hJ{וַיֹּאמֶר מֶה עָשִׂיתָ קוֹל דְּמֵי אָחִיךָ צֹעֲקִים אֵלַי מִן הָאֲדָמָה}
}
{
	\eJ{
		And He said, “What have you done? The sound!
		Your brother’s blood is crying to me from the ground!
	}
}
{
	Creator is mad at creator.
}

\Verse{11}
{
	\hJ{וְעַתָּה אָרוּר אָתָּה מִן הָאֲדָמָה אֲשֶׁר פָּצְתָה אֶת פִּיהָ לָקַחַת אֶת דְּמֵי אָחִיךָ מִיָּדֶךָ}
}
{
	\eJ{
		And now you’re cursed from the ground that opened its mouth
		to take your brother’s blood from your hand.
	}
}
{
	So now he's cursed from the \paleo{אֲדָמָה} (admh) too.

	Like father like son.
}

\Verse{12}
{
	\hJ{כִּי תַעֲבֹד אֶת הָאֲדָמָה לֹא תֹסֵף תֵּת כֹּחָהּ לָךְ נָע וָנָד תִּהְיֶה בָאָרֶץ}
}
{
	\eJ{
		When you work the ground it won’t continue to give its potency to you.
		You’ll be a roamer and rover in the earth.”
	}
}
{
	So he curses him to wander around.
}

\Verse{13}
{
	\hJ{וַיֹּאמֶר קַיִן אֶל יְהֹוָה גָּדוֹל עֲוֺנִי מִנְּשֹׂא}
}
{
	\eJ{
		And Cain said to YHWH, “My crime is greater than I can bear.
	}
}
{
	So Cain regrets it.
}

\Verse{14}
{
	\hJ{הֵן גֵּרַשְׁתָּ אֹתִי הַיּוֹם מֵעַל פְּנֵי הָאֲדָמָה וּמִפָּנֶיךָ אֶסָּתֵר וְהָיִיתִי נָע וָנָד בָּאָרֶץ וְהָיָה כל מֹצְאִי יַהַרְגֵנִי}
}
{
	\eJ{
		Here, you’ve expelled me from the face of the ground today
		and I’ll be hidden from your presence,
		and I’ll be a roamer and rover in the earth,
		and anyone who finds me will kill me.”
	}
}
{
	He's like ``Anyone who finds me will kill me.''
}

\Verse{15}
{
	\hJ{וַיֹּאמֶר לוֹ יְהֹוָה לָכֵן כּל הֹרֵג קַיִן שִׁבְעָתַיִם יֻקָּם וַיָּשֶׂם יְהֹוָה לְקַיִן אוֹת לְבִלְתִּי הַכּוֹת אֹתוֹ כּל מֹצְאוֹ}
}
{
	\eJ{
		And YHWH said to him,
		“Therefore: anyone who kills Cain,
		he’ll be avenged sevenfold.”
		And YHWH set a sign for Cain
		so that anyone who finds him would not strike him.
	}
}
{
	So YHWH says ``Nobody kill this guy.''

	He says if you kill Cain, he'll you like, seven times, or something.
	
	Nobody knows what the sign is.
}

\Verse{16}
{
	\hJ{וַיֵּצֵא קַיִן מִלִּפְנֵי יְהֹוָה וַיֵּשֶׁב בְּאֶרֶץ נוֹד קִדְמַת עֵדֶן}
}
{
	\eJ{
		And Cain went out from YHWH’s presence
		and lived in the land of Nod, east of Eden.
	}
}
{
	So he goes to the land of Wandering.

	It's apparently east of Eden.

	Now back in Gen. 3:24, we learned that the exit to Eden is on the east.

	So it's not entirely clear he went anywhere.

	\aC{
		So Cain goes to Iran.

		Iran is East of Eden.

		No really, I promise.

		Source: Tigris and Euphrates from Genesis 2, then go East.
	}
}

\Verse{17}
{
	\hJ{וַ}
	\hJ{\hlC{יָדַע}}
	\hJ{\hlB{קַיִן}}
	\hJ{אֶת אִשְׁתּוֹ וַתַּהַר וַתֵּלֶד אֶת חֲנוֹךְ וַיְהִי בֹּנֶה עִיר וַיִּקְרָא שֵׁם הָעִיר כְּשֵׁם בְּנוֹ חֲנוֹךְ}
}
{
	\eJ{
		And \hlB{Cain} \hlC{knew} his wife,
		and she became pregnant and gave birth to Enoch.
		And he was a builder of a city,
		and he called the name of the city
		like the name of his son: Enoch.
	}
}
{
	Anyways, Cain has sex with his wife.%
	\footnote{Readers of the bible --- religious and unreligious, scholars and laymen, modern and ancient --- are in universal agreement that this verse raises the highly important theological question, ``Where did she come from?!''}

	And she apparently (1) gets pregnant, and (2) exists.

	The verbs are in the usual order this time.

	That's a bit more evidence that the grammar in Gen 4:1 was intentional.

	Anyways, Cain makes a kid and names it Enoch.

	And he also makes a city, and names it Enoch.
}

\Verse{18}
{
	\hJ{וַיִּוָּלֵד לַחֲנוֹךְ אֶת עִירָד וְעִירָד יָלַד אֶת מְחוּיָאֵל וּמְחִויאֵל יָלַד אֶת מְתוּשָׁאֵל וּמְתוּשָׁאֵל יָלַד אֶת לָמֶךְ}
}
{
	\eJ{
		And Irad was born to Enoch,
		and Irad fathered Mehuya-el,
		and Mehuya-el fathered Metusha-el,
		and Metusha-el fathered Lamech.
	}
}
{
	Then Enoch has a kid and names it ``city.''
	
	\Def{Irad}[\heb{עִירָד}; \textit{ʿÎrāḏ}]{
		Cain's grandson and the son of Enoch (Gen.~4:18). The name is
		probably related to \heb{עִיר} (\textit{ʿîr}, ``city''), though its
		etymology is uncertain.
	}
	
	\Def{City}[\heb{עִיר}; \textit{ʿîr}]{
		City, town, settlement. The standard Hebrew word for an inhabited,
		often fortified, urban center.
	}
	
	Then there are some people with the word ``El'' in their names.
	
	\Def{Mehuja-el}[\heb{מְחוּיָאֵל}; \textit{Məḥûyāʾēl}]{
		Son of Irad. An early theophoric name containing \heb{אֵל} (El).
		Its precise meaning is uncertain, with proposed translations including
		``struck by El'' and ``given life by El.''
	}
	
	\Def{Methusha-el}[\heb{מְתוּשָׁאֵל}; \textit{Məṯûšāʾēl}]{
		Son of Mehujael. Another early theophoric name containing \heb{אֵל}.
		Its exact meaning is uncertain, though it is commonly understood
		as something like ``man of El'' or ``one who belongs to El.''
	}
	
	Lots going on here.
	
	\Def{El}[\heb{אֵל}; \textit{ʾēl}]{
		The name of the chief deity of the Canaanite pantheon, and also the
		generic Northwest Semitic word for ``god.'' Names containing
		\heb{אֵל} were common throughout Canaan long before the worship of
		YHWH became dominant in Israel, and remained common in Israelite
		names such as \textit{Samuel}, \textit{Ezekiel}, \textit{Daniel},
		and especially \textit{Israel} (``El strives'' or ``May El strive'').
	}

	Then there's this guy named LMK (Lamech).

	\Def{Lamech}[\heb{לֶמֶךְ}; \textit{Lemeḵ}]{
		The meaning of his name is unknown; no convincing Hebrew or
		Northwest Semitic etymology has been identified.
	}

}

\Verse{19}
{
	\hJ{וַיִּקַּח לוֹ לֶמֶךְ שְׁתֵּי נָשִׁים שֵׁם הָאַחַת עָדָה וְשֵׁם הַשֵּׁנִית צִלָּה}
}
{
	\eJ{
		And Lamech took two wives.
		The one’s name was Adah,
		and the second’s name was Zillah.
	}
}
{
	So Lamech invents polygamy.
}

\Verse{20}
{
	\hJ{וַתֵּלֶד עָדָה אֶת יָבָל הוּא הָיָה אֲבִי יֹשֵׁב אֹהֶל וּמִקְנֶה}
}
{
	\eJ{
		And Adah gave birth to Jabal.
		He was father of tent-dweller and cattleman.
	}
}
{
	Nobody knows why Jabal.
}

\Verse{21}
{
	\hJ{וְשֵׁם אָחִיו יוּבָל הוּא הָיָה אֲבִי כּל תֹּפֵשׂ כִּנּוֹר וְעוּגָב}
}
{
	\eJ{
		And his brother’s name was Jubal.
		He was father of every player of lyre and pipe.
	}
}
{
	Nobody knows why Jubal.

	\Def{Byblos}[Greek: Βύβλος; also known as \textit{Jbail, Jebeil, Jbeil, Jubayl}; Arabic \arb{جُبَيْل} \textit{Jbeil, Jubayl}]{
		An ancient city in the Keserwan--Jbeil Governorate of Lebanon.
		The area is believed to have been first settled between 8800 and
		7000~BC and continuously inhabited since 5000~BC. During its
		history, Byblos was part of numerous cultures including Canaanite,
		Egyptian, Phoenician, Assyrian, Persian, Hellenistic, Roman,
		Genoese, Mamluk, and Ottoman. Urbanisation is thought to have
		begun during the third millennium BC when it developed into a city,
		making it one of the oldest cities in the world, if not the oldest.
	}
}

\Verse{22}
{
	\hJ{וְצִלָּה גַם הִוא יָלְדָה אֶת תּוּבַל קַיִן לֹטֵשׁ כּל חֹרֵשׁ נְחֹשֶׁת וּבַרְזֶל וַאֲחוֹת תּוּבַל קַיִן נַעֲמָה}
}
{
	\eJ{
		And Zillah, too, gave birth to Tubal-Cain,
		forger of every implement of bronze and iron.
		And Tubal-Cain’s sister was Naamah.
	}
}
{
	Naamah is also Rehoboam's mom.
	
	\aB{Who's Re---}
	
	Forget it. We'll get there eventually.
}

\Verse{23}
{
	\hJ{וַיֹּאמֶר לֶמֶךְ לְנָשָׁיו עָדָה וְצִלָּה שְׁמַעַן קוֹלִי נְשֵׁי לֶמֶךְ הַאְזֵנָּה אִמְרָתִי כִּי אִישׁ הָרַגְתִּי לְפִצְעִי וְיֶלֶד לְחַבֻּרָתִי}
}
{
	\eJ{
		And Lamech said to his wives:
		Adah and Zillah, listen to my voice,
		Wives of Lamech, hear what I say.
		For I’ve killed a man for a wound to me
		And a boy for a hurt to me,
	}
}
{
	So Lamech's like:

	\begin{quotation}
	\aC{
%		Then Lamech spoke; his haughty bosom glowed,\\
%		And thus before his wondering consorts showed:
%
%		``Adah and Zillah, hear your lord's command;\\
%		Attend the words that fall from Lamech's hand.
%
		A man I slew who dared my peace molest;\\
		A youth I struck who rose against my breast.

		The wound he gave was slight, the vengeance great;\\
		Such is the portion rebels earn from fate.
%
%		Let weaker souls endure the threatening blow,\\
%		Or nurse their wrongs and let their anger slow:
%
%		If Cain, though guilty, found such guarding care\\
%		That sevenfold vengeance waited for him there,
%
%		Then far more dreadful shall my sentence prove,\\
%		And Heaven itself my injured right approve.
%
%		Not seven alone shall answer for the deed,\\
%		But seventy times and seven shall justice heed.
%
%		Let all who hear my name remember well:\\
%		I strike but once, yet hosts shall for it fall.''
	}
	\end{quotation}

	\hphantom{So Lamech's like:} ``Look I'm Cain twice.''
}

\Verse{24}
{
	\hJ{כִּי שִׁבְעָתַיִם יֻקַּם קָיִן וְלֶמֶךְ שִׁבְעִים וְשִׁבְעָה}
}
{
	\eJ{
		For Cain will be avenged sevenfold
		And Lamech seventy-seven.
	}
}
{
	Two wives.
	
	Two murders.
	
	Two sevens.
	
	Lamech out.
}

\Verse{25}
{
	\hR{וַיֵּדַע אָדָם עוֹד אֶת אִשְׁתּוֹ וַתֵּלֶד בֵּן וַתִּקְרָא אֶת שְׁמוֹ שֵׁת כִּי שָׁת לִי אֱלֹהִים זֶרַע אַחֵר תַּחַת הֶבֶל כִּי הֲרָגוֹ קָיִן}
}
{
	\eR{
		And Adam knew his wife again,
		and she gave birth to a son,
		and she called his name Seth
		“because {\God} put another seed for me
		in place of Abel because Cain killed him.”
	}
}
{
	Now the Redactor shows up.

	He needs to solve a problem we haven't gotten to yet.

	Stick around for a few more sentences and that'll make sense.
}

\Verse{26}
{
	\hR{וּלְשֵׁת גַּם הוּא יֻלַּד בֵּן וַיִּקְרָא אֶת שְׁמוֹ אֱנוֹשׁ}
	\hJ{אָז הוּחַל לִקְרֹא בְּשֵׁם יְהֹוָה}
}
{
	\eR{
		And a son was born to Seth, him as well,
		and he called his name Enosh.
	}
	\eJ{
		Then it was begun to invoke the name YHWH.
	}
}
{
	Then Seth has a kid whose name is suspiciously similar to Enoch.
	
	Then the chapter ends with a strange comment that seems to say

	This was when people started to use the name YHWH.
}
